% !TeX program = uplatex
\documentclass[a4paper,12pt]{jlreq}

% ===== Packages =====
\usepackage{amsmath,amssymb,mathtools,bm}
\usepackage{siunitx}
\sisetup{detect-all,per-mode=symbol}
\usepackage{physics}
\usepackage{mhchem}
\usepackage{booktabs}
\usepackage{graphicx}
\usepackage{hyperref}
\usepackage{ascmac}
\usepackage{xcolor}


\hypersetup{
  colorlinks=true,
  linkcolor=black,
  urlcolor=blue,
  citecolor=black
}
\usepackage{enumitem}
\usepackage{geometry}
\usepackage{fancybox}
\geometry{margin=25mm}

\title{解析力学ゼミ資料 (考える力学13.10〜13.11)}
\author{学籍番号：2225091　西亦 岳雄}

\begin{document}
\maketitle




\section{13.10 ハミルトンの正準方程式}

\subsection*{13.10 ハミルトンの正準方程式の導出}

\subsubsection{全微分}
1 変数関数 $y = f(x)$ の微小変化は
\[
  dy = \frac{df}{dx}\,dx
\]
で与えられる．

$n$ 個の変数 $x_1,\dots,x_n$ の関数 $f(x_1,\dots,x_n)$ の全部微分は
\[
  df = \sum_{k=1}^n 
        \frac{\partial f}{\partial x_k}\,dx_k
\tag{13.128}
\]
と書かれる．

\subsubsection{ルジャンドル変換}
1 変数関数 $f(x)$ に対し
\[
  y = \frac{df}{dx}
\]
で定義される新しい変数 $y$ を導入し，
\[
  g(y) = xy - f(x)
\]
をルジャンドル変換と呼ぶ．



\subsubsection{一般化運動量とハミルトニアンへの変換}
ラグランジアン $L(q_1,\dots,q_f,\dot q_1,\dots,\dot q_f,t)$
が与えられた力学系を考える．
一般化運動量 $p_k$ を
\[
  p_k = \frac{\partial L}{\partial \dot q_k}
  \qquad (k=1,\dots,f)
\tag{13.133}
\]
で定義する．

ルジャンドル変換の考え方を用いて，
$\dot q_k$ を $p_k$ に置き換えた関数
\[
  H(q_1,\dots,q_f,p_1,\dots,p_f,t)
  = \sum_{k=1}^f p_k \dot q_k - L(q,\dot q,t)
\tag{13.138}
\]
を \textbf{ハミルトニアン}（ハミルトン関数）と呼ぶ．

\subsubsection{ハミルトンの正準方程式}
$H$ の全微分を求めると
\[
  dH = \sum_{k=1}^f
       \frac{\partial H}{\partial q_k}\,dq_k
     + \sum_{k=1}^f
       \frac{\partial H}{\partial p_k}\,dp_k
     + \frac{\partial H}{\partial t}\,dt
\tag{13.139}
\]
である．一方，定義式 $H=\sum p_k\dot q_k - L$ を時間微分し，
オイラー・ラグランジュ方程式
\[
  \frac{d}{dt}
  \frac{\partial L}{\partial \dot q_k}
  - \frac{\partial L}{\partial q_k} = 0
\tag{13.103}
\]
を用いて整理すると
\[
  dH = \sum_{k=1}^f \dot q_k\,dp_k
      - \sum_{k=1}^f \dot p_k\,dq_k
      - \frac{\partial L}{\partial t}\,dt
\tag{13.143}
\]
が得られる．
この 2 つの表現を $dq_k,dp_k,dt$ それぞれの係数で比較すると
\[
  \dot q_k = \frac{\partial H}{\partial p_k},\qquad
  \dot p_k = -\frac{\partial H}{\partial q_k}
  \qquad (k=1,\dots,f)
\tag{13.144}
\]
\[
  \frac{\partial H}{\partial t}
  = -\frac{\partial L}{\partial t}
\tag{13.140}
\]
が従う．
特に $L$ が時間に依存しないとき，
$\partial L/\partial t=0$ なので
\[
  \frac{dH}{dt} = \frac{\partial H}{\partial t} = 0
\]
となり，$H$ は保存量である．これ(13.144)が
\textbf{ハミルトンの正準方程式}である．

\subsubsection{ハミルトニアンとエネルギー}
質点系のラグランジアンが
\[
  L = K - U
\]
（$K$ は運動エネルギー，$U$ はポテンシャル・エネルギー）で与えられ，
かつ一般化座標が直交座標と同等の性質をもつ場合には，
一般化運動量と速度の関係は
\[
  p_i = m_i \dot{\bm r}_i
\tag{13.145}
\]
であり，
\[
  \frac{1}{2}\sum_i p_i\cdot\dot{\bm r}_i = K
\tag{13.150}
\]
が成り立つ．したがってハミルトニアンは
\[
  H = \sum_i p_i\cdot\dot{\bm r}_i - L
    = K + U
\tag{13.152}
\]
となり，全力学系の力学的エネルギーに等しい．
$L$ が時間に依存しないとき，$H$ は保存エネルギーを表す．

\subsection*{13.11 ポアソンの括弧}

\subsubsection{定義}
正準変数 $(q_1,\dots,q_f,p_1,\dots,p_f)$ の関数
$F(q,p,t)$ と $G(q,p,t)$ に対し，
\textbf{ポアソンの括弧}を
\[
  [F,G]
  = \sum_{k=1}^f
    \left(
      \frac{\partial F}{\partial q_k}
      \frac{\partial G}{\partial p_k}
      - 
      \frac{\partial F}{\partial p_k}
      \frac{\partial G}{\partial q_k}
    \right)
\tag{13.155}
\]
で定義する．

\subsubsection{時間発展との関係}
$F(q,p,t)$ の時間微分は
ハミルトンの正準方程式を用いると
\[
  \frac{dF}{dt}
  = \sum_{k=1}^f
    \left(
      \frac{\partial F}{\partial q_k}\dot q_k
      +\frac{\partial F}{\partial p_k}\dot p_k
    \right)
  + \frac{\partial F}{\partial t}
\]
\[
  = [F,H] + \frac{\partial F}{\partial t}
\tag{13.156}
\]
と書ける．
特に $F$ が時間に依存しないとき，
\[
  \frac{dF}{dt} = [F,H]
\tag{13.157}
\]
となる．
このとき
\[
  [F,H]=0
\tag{13.158}
\]
であれば $F$ は保存量である．

特に, (13.157)に関して$F = q_k, p_k$であるとき

\[
  \dot{q_k} = [q_k,H],　\dot{p_k} = [p_k,H],　(k = 1,..., n)
  \tag{13.159}
\]
で, これはポアソン括弧を用いた正準方程式である.









 


\begin{thebibliography}{9}

\bibitem{tasaki}

考える力学



\end{thebibliography}

\end{document} 